\documentclass[11pt,a4paper]{article}
\usepackage[T1]{fontenc}
\usepackage[utf8]{inputenc}
\usepackage{lmodern,microtype}
\usepackage{amsmath,amssymb,mathtools,amsthm}
\usepackage{geometry,booktabs,tabularx,hyperref}
\geometry{margin=2.4cm}
\hypersetup{colorlinks=true,linkcolor=black,citecolor=black,urlcolor=blue}
\newtheorem{theorem}{Theorem}[section]
\newtheorem{proposition}[theorem]{Proposition}
\newtheorem{remark}[theorem]{Remark}
\setlength{\parindent}{1.2em}
\setlength{\parskip}{0.15em}

\title{\textbf{Random and Structured Continuous SAT Dynamics}\\[0.3em]
\large LP-Volume Bounds, Horn Chains, Subcritical Hinge Residuals, and a Complexity Firewall}
\author{Thiago Carvalho\\Independent Researcher\\\texttt{thiagocarvalho.dba@gmail.com}}
\date{Draft for analysis -- derived from audited CLG-R v4.0.5 source (September 2026)}

\begin{document}
\maketitle

\begin{abstract}
We collect several results on random and structured 3-SAT instances under continuous projected dynamics, with emphasis on what can and cannot be inferred from low-dimensional trapping phenomena. For the independent-clause model, we prove by Fubini, the Irwin--Hall law and Jensen's inequality that the expected normalized volume of the squared-Hinge zero set is bounded below by $(5/6)^{\lfloor\alpha N\rfloor}$. In the uniform sparse model below the branching threshold $\alpha<1/6$, isolated clauses occur with normalized density $e^{-9\alpha}$. Their three-dimensional projected Hinge dynamics has exact Boolean rounding failure probability $3/32$, yielding a positive residual lower bound $(3/32)e^{-9\alpha}$ in expectation and with asymptotically high probability. We then examine Horn implication chains. Pure chains conserve the coordinate mean, while an identification of their continuous limit with Euclidean isotonic projection remains unproved; the associated Sparre--Andersen expression is therefore only conditional. Adding a positive head fact destroys mean conservation and collapses the zero set to a singleton, falsifying a previously contemplated large trapping basin on this family. General Horn clauses have mixed Jacobian signs, so neither standard cooperative nor standard competitive monotonicity theory settles their dynamics. We conclude with an epistemological firewall: continuous glassy or metastable behavior is not a proxy for Turing complexity, as illustrated by the polynomial-time solvability of 3-XOR-SAT despite glassy continuous landscapes observed in the associated physics literature and experiments. The paper is intended as a rigorous research note separating proven random-geometry statements from open dynamical conjectures.
\end{abstract}

\noindent\textbf{Keywords.} random 3-SAT; Hinge relaxation; Horn-SAT; projected dynamics; LP polytope; XORSAT; spin glasses.

\section{Introduction}
Continuous relaxations of SAT mix three logically distinct questions: geometry of the zero-energy relaxation, dynamics of a particular flow, and computational complexity of the underlying discrete problem. A central objective of this paper is to keep these levels separate.

We begin with a finite-dimensional geometric statement for the clause-wise LP polytope induced by the squared-Hinge relaxation. We then isolate a subcritical random-formula mechanism that creates a positive residual density under Hinge rounding: isolated clauses. The mechanism is completely local and analytically solvable. Next we revisit implication chains and general Horn clauses, recording both rigorous statements and a remaining open identification problem. Finally, we explain why dynamical trapping cannot by itself be used as evidence for a complexity-class separation.

Throughout,
\[
g_c(x)=-\frac12(1+\sigma^{(c)}\cdot x),
\qquad
\Phi_{\mathrm{quad}}(x)=\sum_c[\max(0,g_c(x))]^2,
\]
on $\mathcal X=[-1,1]^N$, with projected flow $\dot x=\Pi_{T_{\mathcal X}(x)}(-\nabla\Phi_{\mathrm{quad}}(x))$. We distinguish the independent-clause model $\mathcal E_{\mathrm{iid}}(N,\alpha)$ from the uniform distinct-clause model $\mathcal E_{\mathrm{unif}}(N,\alpha)$.

\section{A finite-dimensional lower bound on LP-polytope volume}
Let
\[
Z(F)=\{x\in[-1,1]^N:g_c(x)\le0\text{ for every clause }c\}.
\]
\begin{theorem}[Jensen lower bound]
For $F\sim\mathcal E_{\mathrm{iid}}(N,\alpha)$, with $M=\lfloor\alpha N\rfloor$ and $N\ge3$,
\[
\mathbb E_F\mu_{\mathrm{norm}}(Z(F))
\ge\left(\frac56\right)^{\lfloor\alpha N\rfloor}
\ge e^{-\alpha N\ln(6/5)}>0.
\]
Also $\mathbb E\mu_{\mathrm{norm}}(Z)\ge(1/3)^N$, because the universal central box $(-1/3,1/3)^N$ is contained in $Z$ for every formula.
\end{theorem}
\begin{proof}
Fubini and clause independence give
\[
\mathbb E_F\mu_{\mathrm{norm}}(Z)
=\int_{[-1,1]^N}p(x)^M\frac{dx}{2^N},
\]
where $p(x)$ is the probability that one random clause satisfies the continuous LP inequality at $x$. Averaging jointly over $x$ and random literal signs turns the three literal-satisfaction coordinates into i.i.d. $U[0,1]$ variables. Their sum $S_3$ satisfies
\[
\mathbb P(S_3<1)=\frac16,
\]
the volume of the standard three-dimensional simplex, so $\int p(x)dx/2^N=5/6$. Jensen for the convex function $t\mapsto t^M$ yields the claim.
\end{proof}

The theorem is intentionally one-sided. It supplies a nonzero finite-$N$ lower bound but does not determine the asymptotic exponential rate of the true expected volume.

\section{Exact isolated-clause dynamics in the subcritical model}
We next switch to $\mathcal E_{\mathrm{unif}}(N,\alpha)$ and assume $\alpha<1/6$. The clause-exploration branching factor is $6\alpha<1$, so component sizes have an exponential tail and isolated components have positive density.

Consider one isolated signed 3-clause, and transform coordinates by its literal signs so the active halfspace is
\[
S=x_1+x_2+x_3<-1.
\]
Outside the zero polytope all three transformed coordinates move at the same rate. If the initial sum is $S_0<-1$, the trajectory lands on $\partial Z$ at
\[
U_p^*=U_p(0)+\frac{-1-S_0}{3}.
\]
For almost every initial point the landing occurs strictly before any coordinate hits the box boundary.

\begin{proposition}[Exact rounding-failure probability]
For a uniformly initialized isolated clause,
\begin{align}
p_{\mathrm{fail}}
&=\mathbb P(x_0\in Z\text{ and rounding violates})\notag\\
&\quad+\mathbb P(x_0\notin Z\text{ and the landing rounds to a violation})\notag\\
&=\frac1{48}+\frac7{96}=\frac3{32}.
\end{align}
\end{proposition}
The source analysis obtains the two volumes by direct integration in the cube $[-1,1]^3$.

In the sparse uniform ensemble,
\[
\mathbb E|\mathcal C_{\mathrm{iso}}|=\alpha N e^{-9\alpha}(1+O(N^{-1})).
\]
Normalizing by $M\sim\alpha N$ removes the factor $\alpha$.

\begin{theorem}[Subcritical Hinge residual]
For $0<\alpha<1/6$,
\[
\liminf_{N\to\infty}\mathbb E\rho_{\mathrm{quad}}(\alpha)
\ge\frac3{32}e^{-9\alpha}>0.
\]
Moreover, for every $\varepsilon>0$,
\[
\mathbb P\!\left(\rho_{\mathrm{quad}}(\alpha)
\ge\frac3{32}e^{-9\alpha}-\varepsilon\right)\to1.
\]
\end{theorem}
\begin{proof}[Proof sketch]
Conditional on the formula, isolated clauses have disjoint supports and their initialization blocks are independent. Each fails after rounding with probability $3/32$. The isolated-clause count has variance $O(N)$ in the subcritical model; hence its normalized density concentrates by Chebyshev. Combining the concentration of the count with independent local failure variables yields the stated lower bound.
\end{proof}

This result is local: it does not characterize the entire Hinge residual, only a certified contribution from isolated clauses.

\section{Implication chains: a correct conservation law and an open identification}
Consider the chain
\[
x_1\to x_2\to\cdots\to x_K,
\qquad
\Phi_{\mathrm{chain}}(x)=\sum_{k=1}^{K-1}\left[\max\left(0,\frac{x_k-x_{k+1}}2\right)\right]^2.
\]
The internal forces telescope, so
\[
\frac{d}{dt}\sum_{k=1}^Kx_k(t)=0.
\]
Thus pure-chain trajectories remain on affine hyperplanes of fixed mean.

The zero set is the isotonic cone intersected with the cube,
\[
Z=\{x_1\le x_2\le\cdots\le x_K\}.
\]
A tempting conjecture is that the continuous Hinge flow converges to the Euclidean isotonic projection $\Pi_Z(x_0)$. The audited source does not prove this identity, and we retain it as an open analytical problem.

If that identity is assumed, then the first projected coordinate is positive exactly when the relevant prefix-sum conditions are positive, leading to the classical Sparre--Andersen survival expression
\[
\mathbb P(x_1^*>0\mid x^*=\Pi_Z(x_0))
=\frac{\binom{2K}{K}}{4^K}
=\frac1{\sqrt{\pi K}}\left(1-\frac1{8K}+O(K^{-2})\right).
\]
This is therefore conditional evidence, not an unconditional theorem about the continuous flow. A submission version should add a primary bibliographic citation for the classical Sparre--Andersen result.

\section{A positive head fact destroys the proposed chain trap}
Add a positive unit fact through
\[
h(x_1)=\left[\max\left(0,\frac{1-x_1}{2}\right)\right]^2.
\]
For $x_1<1$,
\[
-\partial_{x_1}h=\frac{1-x_1}{2}>0,
\]
so the mean conservation law is broken. More decisively, the zero set of the resulting convex squared-Hinge objective collapses to
\[
Z=\{(+1,\ldots,+1)\}.
\]
Projected convex-gradient-flow convergence then forces every trajectory to this singleton. Hence a previously contemplated $1-o(1)$ spurious trapping basin on implication chains with a positive head fact is false. This negative result identifies a family that cannot witness the desired Hinge-versus-multilinear separation.

\section{Horn clauses and Jacobian sign patterns}
For a purely negative clause
\[
(\neg x_i\lor\neg x_j\lor\neg x_k),
\qquad
P_c(x)=\frac{(1+x_i)(1+x_j)(1+x_k)}8,
\]
we have
\[
\partial_{ij}^2P_c=\frac{1+x_k}{8}\ge0,
\]
so the off-diagonal Jacobian entries of the unprojected multilinear vector field $-\nabla\Phi_{\mathrm{mult}}$ satisfy $J_{ij}\le0$. Purely negative families therefore have a weakly competitive sign pattern.

General Horn clauses do not preserve one sign. For example, for
\[
(\neg x_1\lor\neg x_2\lor x_3)
\]
one obtains
\[
J_{12}=-\frac{1-x_3}{8}\le0,
\qquad
J_{13}=\frac{1+x_2}{8}\ge0,
\qquad
J_{23}=\frac{1+x_1}{8}\ge0.
\]
Thus general Horn 3-SAT is neither automatically cooperative nor automatically competitive with respect to the standard positive orthant order. Standard monotone-system theorems do not by themselves settle the global dynamics.

\begin{remark}[Open problem]
The exact spurious-basin mass for general Horn DAGs remains open in the audited source. The pure implication chain with a positive head fact cannot solve the problem because its zero set is a singleton.
\end{remark}

\section{Relation to the multilinear obstruction}
The isolated-clause Hinge lower bound above and the six-clause M6 multilinear obstruction answer different questions. The Hinge mechanism is driven by the large fractional zero polytope and already appears on an isolated single clause. The multilinear mechanism has no interior plateau; its certified positive residual instead arises from a six-clause boundary equilibrium family. In the uniform random ensemble, the audited M6 theorem gives
\[
\liminf_{N\to\infty}\mathbb E\rho_{N,\mathrm{mult}}(\alpha)
\ge\frac{729}{2^{59}}\alpha^5e^{-39\alpha}>0
\]
for every fixed $\alpha>0$. Thus positivity of either residual is not the remaining comparison problem; the unresolved random-ensemble question is whether one residual is strictly larger than the other in the relevant density regime.

\section{A complexity firewall: 3-XOR-SAT}
Continuous dynamics should not be used as a surrogate definition of computational hardness. The audited CLG-R source emphasizes 3-XOR-SAT as a firewall example. A system of XOR clauses is a linear system over $\mathbb F_2$ and is solvable in deterministic polynomial time by Gaussian elimination. Yet the associated continuous multilinear landscapes can exhibit glassy trapping in experiments, and diluted $p$-spin/XORSAT models have a developed metastability literature \cite{mezard2003two,ricci2010being}.

The logical conclusion is modest but important:
\begin{quote}
Failure or success of a particular continuous gradient flow neither implies nor is implied by Turing solvability of the underlying discrete problem.
\end{quote}
Consequently none of the geometric or dynamical separation phenomena discussed here constitute evidence for $P\ne NP$ or $P=NP$.

\section{Research agenda}
The source material supports several sharply formulated next questions:
\begin{enumerate}
\item prove or refute the isotonic-projection identity for continuous Hinge implication-chain flow;
\item characterize global dynamics of general Horn DAGs with mixed Jacobian signs;
\item extend random residual comparisons beyond isolated local motifs;
\item establish motif-frequency theorems for planted ensembles rather than transferring results from the unplanted uniform model;
\item determine the asymptotic LP-polytope volume rate, since the Jensen theorem is only a lower bound;
\item compare the Hinge and multilinear residual densities in the random clustering regime without conflating positivity with strict separation.
\end{enumerate}

\section{Conclusion}
The random and structured results in this note separate several mechanisms that can otherwise be conflated. The Hinge LP polytope has a provably positive finite-dimensional expected-volume lower bound. Isolated clauses alone certify a positive subcritical Hinge residual. Pure Horn implication chains obey an exact conservation law, but their proposed isotonic limit remains unproved, and a positive unit fact eliminates the suspected large trapping basin. General Horn clauses have mixed interaction signs. Finally, 3-XOR-SAT shows why continuous dynamical difficulty and discrete computational complexity must remain conceptually separate.

\begin{thebibliography}{9}
\bibitem{goemans1995improved} M. X. Goemans and D. P. Williamson, "Improved approximation algorithms for maximum cut and satisfiability problems using semidefinite programming", \emph{Journal of the ACM}, 42(6):1115--1145, 1995.
\bibitem{krzakala2007gibbs} F. Krzakala, A. Montanari, F. Ricci-Tersenghi, G. Semerjian, and L. Zdeborov\'a, "Gibbs states and the set of solutions of random constraint satisfaction problems", \emph{Proceedings of the National Academy of Sciences}, 104(25):10318--10323, 2007.
\bibitem{mezard2003two} M. M\'ezard, F. Ricci-Tersenghi, and R. Zecchina, "Two solutions to diluted p-spin models and XORSAT problems", \emph{Journal of Statistical Physics}, 111(3):505--533, 2003.
\bibitem{ricci2010being} F. Ricci-Tersenghi, "Being glassy without being hard to solve", \emph{Science}, 330(6011):1639--1640, 2010.
\bibitem{brogliato2006equivalence} B. Brogliato, A. Daniilidis, C. Lemar\'echal, and J. Malick, "Equivalence and complementarity issues in projected dynamical systems and differential inclusions", \emph{Mathematical Programming}, 107(3):317--335, 2006.
\bibitem{nagurney1996projected} A. Nagurney and D. Zhang, \emph{Projected Dynamical Systems and Variational Inequalities with Applications}. Springer, 1996.
\end{thebibliography}
\end{document}
